\section{Physical Results \realmark}
\label{sec:realresults}

\subsection{What was flown, and what was recorded}

Nine runs were flown on a single day, all with the committed planner.
Run~5 is excluded by the criterion described below, leaving eight
admissible runs across three geometries: \textsf{K0} with the anchor on
the approach line ($n=4$), \textsf{K1} offset \SI{0.35}{\metre} ($n=2$),
and \textsf{K1M}, the mirror of \textsf{K1} ($n=2$). \textsf{K1M} was
invented during the session and was never pre-registered; its purpose
and its limits are stated in Section~\ref{sec:sidechoice}.

Each run recorded the command stream at \SI{19}{\hertz} and an annotated
onboard video. \emph{No external tracking was recorded during any run.}
The overhead camera was used to measure and verify the start pose before
each release and was deliberately not opened inside the run process,
because doing so had killed several runs outright. The consequence is
permanent and applies to every physical run in this paper: trajectory,
path length, lateral excursion, forward progress, true range to the
anchor and true minimum clearance do not exist and cannot be
reconstructed.

We could have integrated the commanded velocity and produced all six.
We do not, here or anywhere else, because integrating an open-loop model
of a saturating, deadbanded vehicle in moving water would manufacture
precisely the quantities under test. Where a quantity was not observed,
the tables below say so.

\subsection{Admissibility}

Some recordings contain more command samples than the \SI{19}{\hertz}
publication grid holds, because orphaned planner processes from earlier
runs survived a name-based process sweep and published to the same
topic. Where the surplus samples are exact duplicates they carry no
information and collapsing them onto the grid is resampling; where they
are not, they are genuine divergence between publishers and the run
cannot be repaired. The admission test is measured rather than chosen:
the threshold is the largest within-cell command spread among recordings
made \emph{at} the grid rate, which have a single publisher by
construction (\SI{0.0107}{\metre\per\second}). Eight runs pass; run~5,
at \SI{0.0349}{\metre\per\second}, does not, and is excluded while its
recording is retained.

\subsection{Command-domain results}

Table~\ref{tab:real} reports every admissible run individually. With
eight runs and two per off-axis geometry, we report runs and medians and
avoid inferential statistics entirely; an effect size computed on $n=2$
would be decoration.

\begin{table}[t]
\centering
\caption{The eight admissible physical runs \realmark. \emph{det}:
fraction of frames in which the detector reported the obstacle;
\emph{surge}: mean commanded surge; \emph{lat}: fraction of samples with
a lateral command; \emph{rev}: sign reversals; \emph{side} and $t_c$:
side and time of the first sustained lateral commitment
($>\SI{0.05}{\metre\per\second}$ held \SI{1.0}{\second}). Entries marked
\emph{cens.}\ began already committed, so the time is a bound.
Commitment \emph{distance} is absent from this table because it was
never observed.}
\label{tab:real}
\begin{tabular}{rlrrrrrlr}
\toprule
run & geo & Hz & det & surge & lat & rev & side & $t_c$ (\si{\second}) \\
\midrule
1 & K0  & 18.9 & 0.79 & 0.149 & 0.32 & 1 & right & 8.84 \\
2 & K0  & 18.8 & 0.89 & 0.113 & 0.53 & 2 & right & 5.16 \\
3 & K0  & 36.2 & 0.88 & 0.133 & 0.24 & 0 & right & cens. \\
4 & K0  & 55.6 & 0.97 & 0.113 & 0.53 & 2 & right & cens. \\
6 & K1  & 18.7 & 0.94 & 0.105 & 0.64 & 1 & right & 6.05 \\
7 & K1  & 18.7 & 0.85 & 0.123 & 0.38 & 1 & left  & 10.53 \\
8 & K1M & 18.6 & 0.84 & 0.120 & 0.43 & 1 & left  & 6.58 \\
9 & K1M & 18.7 & 0.71 & 0.110 & 0.57 & 3 & left  & 8.53 \\
\bottomrule
\end{tabular}
\end{table}

Three observations. The vehicle detected the anchor in
\SIrange{71}{97}{\percent} of frames, and no collision was recorded in
an admissible run; commanded surge is zero in less than
\SI{1}{\percent} of samples throughout. Median commanded surge is
\SI{0.123}, \SI{0.114} and \SI{0.115}{\metre\per\second} in
\textsf{K0}, \textsf{K1} and \textsf{K1M} -- propulsion does not depend
on geometry, as it should not. And two of the four \textsf{K0} runs
are left-censored: they were already commanding laterally at the first
recorded sample, which follows directly from starting
\SI{6}{\centi\metre} outside the engagement range.

\subsection{Side choice follows the obstacle}
\label{sec:sidechoice}

Through the first seven runs the vehicle chose the same side six times.
Two explanations fit that equally well -- the planner selects its side
from where it sees the obstacle, or the vehicle simply always goes right
for a mechanical, perceptual or tie-breaking reason -- and neither
\textsf{K0} nor \textsf{K1} can separate them, because \textsf{K0} puts
the anchor dead ahead and \textsf{K1}'s offset lies on the side a
rightward bias would already favour. Mirroring the offset separates
them in a single run, so the geometry was mirrored and two runs flown.

Both went left, against four of four going right in \textsf{K0}. Read
with its sample size, this is consistent with side choice following the
obstacle and inconsistent with a fixed bias, and it is the strongest
statement two runs can carry. \textsf{K1M} was not pre-registered, was
designed after seeing six same-side outcomes, and has $n=2$; we treat
it as an exploratory observation that motivated the simulated study in
Section~\ref{sec:diagnostic}, not as a confirmatory result.

\subsection{The recovered perception stream}

The detector drew its own output on every frame it processed -- a green
box for an accepted detection, orange for a rejected one, nothing when
it saw nothing -- so the observation stream that reached the planner can
be read back exactly from the recordings, for the six runs that have
video. This is recovery, not re-detection: no detector is re-run, and
the numbers are what the planner was actually given. Reading the box
height through Equation~\ref{eq:range} recovers the range the planner
\emph{believed}, which is the quantity that exists on both sides of the
sim-real comparison. The recovery's own measurement error, from the
four-pixel stroke width, has a median of \SI{0.02}{\metre}, two orders
below the estimator's own bias.

The recovered stream shows what the range window of
Section~\ref{sec:deployment} implies. Accepted detections occupy
\SIrange{71}{91}{\percent} of frames and their median believed range is
\SIrange{0.94}{1.51}{\metre}; between \SI{0} and \SI{28}{\percent} of
them sit exactly at the detector's upper size gate, where the reported
range saturates at \SI{0.30}{\metre}. We report that saturation fraction
rather than a minimum range, because the minimum is the gate and not the
anchor.
